\section{Discussion}
\label{sec:discussion}

\rev{
\myparagraph{Extensibility to new models}
Despite the plethora of model sizes from 2B parameters ~\cite{phi2,zhang2022opt} to 176B parameters~\cite{scao2022bloom} or more~\cite{chatgptapi}, all modern transformer-based generative LLMs have the distinct prompt processing and token generation phases. 
Similarly, even modifications and flavors like Mixture-of-Experts (MoEs) have these phases.
Since Splitwise is built solely by exploiting these phases, it is applicable to all of the current and upcoming LLMs, as long as the auto-regressive nature of the workload requires these two phases.
Note that as shown in \Cref{subsec:exp-workload-changes}, clusters provisioned with Splitwise for one model can also efficiently serve other models. 
%Supporting MoEs (\eg, Mixtral~\cite{jiang2024mixtral}) may require extending our KV-cache transfer implementation. Smaller models (\eg, Phi-2~\cite{phi2}) benefit from multiple instances per server, which we include in our evaluation of Llama2-70B. 
}

\rev{
\myparagraph{Alternative compute hardware}
In this work, we use NVIDIA H100 and A100 GPUs since they are commonly used for LLM inference in datacenters today~\cite{nvidia_hopperincloud}.
Smaller datacenter GPUs like NVIDIA T4 lack enough memory to run modern LLMs efficiently.
In general, our methodology is applicable to any hardware (including CPUs, FPGAs, ASICs~\cite{davidpatterson2019_Domain}) that aligns with the computational requirements of prompt and token phases.
Our characterization suggests that prompt phases need high compute capability and memory bandwidth with low memory capacity, whereas token phases need moderate compute capability with high memory capacity and bandwidth.
Thus, GPUs like AMD MI-250~\cite{mi250} and CPUs like Intel Sapphire-Rapids (with HBM)~\cite{sprhbm} could be effective token machines.
Since we do not have access to such hardware and/or optimized LLM implementations, we leave this to future work.
}
%\rev{
%Since we do not have access to such hardware, we could estimate the performance impact by scaling our profiled results. However, we would be unable to validate our estimates.
%}

%\myparagraphnodot{What would optimized hardware look like?}
%Splitwise enables this hardware design space exploration for each phase independently.

\myparagraph{Interconnect between prompt and token machines}
In this work, we assume Infiniband connection between the prompt and token machines in all the designs (albeit, lower bandwidth when A100s were involved).
Although this is common for all homogenous machines, Splitwise-HA is not be readily available with an Infiniband connection between H100s and A100s, even though technically feasible.
The alternative could be HPC clouds, with Infiniband connections through the CPU~\cite{azurehpcvm}, or Ethernet, using RoCE~\cite{roce}.
Given our optimized KV-cache transfer that helps reduce critical latency, an interconnect with $10\times$ lower bandwidth would likely still be beneficial.
To further reduce our bandwidth utilization, we could also compress the KV-cache before transferring it across the network~\cite{liu2023deja}.

\myparagraph{Heterogeneous prompt/token machines}
Although Splitwise is robust to varied models and input traces, we recognize that fragmenting a data center with different types of GPUs (\eg, Splitwise-HA) may bring its own challenges for the CSP.

\myparagraph{Conversation back and forth}
Chat APIs for LLMs today require the user to send the complete context of the conversation so far~\cite{chatgptapi}.
However, in the future, services may have enough GPU capacity to cache the context and avoid recomputation.
This could sway the memory utilization pattern of the prompt phase from our characterization.
Furthermore, it may require transferring the KV-cache back to a prompt machine to be ready for the next conversation request.
